% Translations and transliteration.
% Compile with: xelatex quran-translations.tex
%
% trans={...} adds an lt/en/de/fr/fa twin of every macro, e.g. \quranayahen.
% Companion packages add more languages: quran-ur, quran-bn, quran-de,
% quran-en, quran-es, quran-id. Loading many at once may exhaust TeX memory
% (see "TeX capacity exceeded" in the manual).
\documentclass{article}
\usepackage[margin=1.8cm]{geometry}
\usepackage{multicol}
\usepackage{forloop}
\usepackage{arabxetex}
\newfontfamily\arabicfont[Script=Arabic, Scale=1.1]{Amiri}
\newfontfamily\bengalifont{Noto Sans Bengali}
\usepackage[uthmani, trans={lt, en, de, fr, fa}]{quran}
\usepackage{quran-ur}
\usepackage{quran-bn}

\parindent=0pt

\begin{document}

\section*{Ayah by ayah: Surah \surahname[103] with English}
% Interleave the original and the translation, one ayah at a time.
\newcounter{ayah}
\forloop{ayah}{1}{\value{ayah} < 4}{%
  \begin{arab}[utf]\quranayah[103][\value{ayah}]\end{arab}
  {\small\itshape\quranayahen[103][\value{ayah}]}\par\medskip}

\section*{Surah \surahname[1] in every language}
\begin{multicols}{2}
  \subsection*{Arabic}
  \begin{arab}[utf]\quransurah[1]\end{arab}

  \subsection*{Transliteration}
  \quransurahlt[1]

  \subsection*{English}
  \quransurahen[1]

  \subsection*{German}
  \quransurahde[1]

  \subsection*{French}
  \quransurahfr[1]

  \subsection*{Persian}
  \begin{arab}[utf]\quransurahfa[1]\end{arab}

  \subsection*{Urdu}
  \begin{arab}[utf]\quransurahur[1]\end{arab}

  \subsection*{Bengali}
  {\bengalifont\quransurahbn[1]}
\end{multicols}

\end{document}
